\section{Census input and tensor conventions}\label{sec:enumeration}
\subsection{One authoritative counting source}
All counting data are taken from the public repository
\href{https://github.com/sebastienpalcoux/fusion-ring-census}{\emph{Fusion Ring Census}}
\cite{CensusRepo}, pinned to revision
\begin{center}\filename{2a340871242879665fa7ba2b01fa15b8bdeaada3}.\end{center}
The source files are \filename{results/census.json} and
\filename{results/rank6/tables.txt.gz}. They declare a complete, verified
rank-six census through multiplicity nine. Its $25421$ tensors include
all $16612$ classes through multiplicity eight and $8809$ classes of exact
multiplicity nine. The exact-multiplicity definition, inclusion of the
unit in the rank, and inclusion of every duality type are unchanged.

The paper trusts this census as its counting input. It supplies neither
another census generator nor a claim to have independently enumerated
the same domain. The repository contains the enumeration algorithms,
coverage arguments, completion records, and tensor checks supporting
Theorem~\ref{thm:census}. A valid and duplicate-free tensor list alone
would not establish exhaustive coverage; that claim is supplied by the
source census, not by the import operation below.

\subsection{Frozen data and reproducible identifiers}
A copy of the source tensor export is retained with the classification
code, so a classification check has no live-network dependency. Its
compressed SHA-256 is
\begin{center}\small
\filename{0e7552e17620452f237e94751c40cd93b0757a466168629db8cc1bf0b9451d00},
\end{center}
and its uncompressed SHA-256 is
\begin{center}\small
\filename{be377c8114d4c0a532b79850efd588090c1dfb5ae3ffc847ecf3eadb97f88ab3}.
\end{center}
These hashes identify the input bytes; they are not substitutes for the
source's completeness argument or for categorification certificates.

Every ring is represented by its full ordered tensor
$N=(N_{ij}^{k})_{0\leq i,j,k<6}$, with $b_0=\one$. The multiplicity is
$\max_{i,j,k}N_{ij}^{k}$. Noncommutative inputs are retained without
symmetrization. A based-ring isomorphism is a unit-fixing basis
permutation preserving every structure constant. Consequently,
comparison with a model uses the complete tensor, not only dimensions,
a fusion graph, or a numerical fingerprint.

The identifier $\RR{m}{j}$ is local to this paper's frozen tensor
registry. Representatives and identifiers through multiplicity eight
are retained exactly, including the fields used by the earlier certificates.
The new exact-multiplicity-nine tensors are placed in
row-major lexicographic order using the exported representatives.
This order need not give the lexicographically smallest representative
under all basis permutations. Import matching records the correspondence
between the source representatives and the registry. Every mathematical
criterion receives the tensor rather than the identifier, so relabelling
cannot alter its conclusion.

The import checks all $25421$ tensors for the fusion-ring axioms and
examines all $120$ unit-fixing permutations of each tensor, for a total of
$3050520$ canonicalization comparisons. It finds no repeated based-ring
class. The explicit source-line and basis-permutation correspondence is
retained in the supplement. These checks establish validity, distinctness,
and agreement with the earlier registry; exhaustive enumeration remains
the separate input described above.

\subsection{Counting, exclusions and constructions}
The source census answers a finite counting question about based rings.
A categorification exclusion answers a different question about complex
fusion categories. We therefore keep the census fixed while applying
necessary criteria and matching models. The resulting ledger covers
every imported tensor exactly once; accepted negative certificates,
positive constructions and unresolved records are disjoint.

The exact prefix $39,154,384,872$ agrees with the complete rank-six range
of \cite{VS}. We use no external categorifiability flag as a proof.
The classification code may reconstruct arithmetic, tensor or
isomorphism data required by a criterion, but it does not repeat the
census search. The counting proposal in \filename{oeis/} derives directly
from the same source counts and has no dependence on whether a
categorification has been found.
